%%
% This is the Eisvogel pandoc LaTeX template.
%
% For usage information and examples visit the official GitHub page:
% https://github.com/Wandmalfarbe/pandoc-latex-template
%%
% Options for packages loaded elsewhere
\PassOptionsToPackage{unicode}{hyperref}
\PassOptionsToPackage{hyphens}{url}
\PassOptionsToPackage{dvipsnames,svgnames,x11names,table}{xcolor}
\documentclass[
  paper=a4,
  ,captions=tableheading
]{scrartcl}
\usepackage{xcolor}
\usepackage[margin=2.5cm,includehead=true,includefoot=true,centering,]{geometry}
\usepackage{amsmath,amssymb}


% add backlinks to footnote references, cf. https://tex.stackexchange.com/questions/302266/make-footnote-clickable-both-ways
\usepackage{footnotebackref}
\setcounter{secnumdepth}{-\maxdimen} % remove section numbering
\usepackage{iftex}
\ifPDFTeX
  \usepackage[T1]{fontenc}
  \usepackage[utf8]{inputenc}
  \usepackage{textcomp} % provide euro and other symbols
\else % if luatex or xetex
  \usepackage{unicode-math} % this also loads fontspec
  \defaultfontfeatures{Scale=MatchLowercase}
  \defaultfontfeatures[\rmfamily]{Ligatures=TeX,Scale=1}
\fi
\usepackage{lmodern}
\ifPDFTeX\else
  % xetex/luatex font selection
\fi
% Use upquote if available, for straight quotes in verbatim environments
\IfFileExists{upquote.sty}{\usepackage{upquote}}{}
\IfFileExists{microtype.sty}{% use microtype if available
  \usepackage[]{microtype}
  \UseMicrotypeSet[protrusion]{basicmath} % disable protrusion for tt fonts
}{}
% Use setspace anyway because we change the default line spacing.
% The spacing is changed early to affect the titlepage and the TOC.
\usepackage{setspace}
\setstretch{1.2}
\makeatletter
\@ifundefined{KOMAClassName}{% if non-KOMA class
  \IfFileExists{parskip.sty}{%
    \usepackage{parskip}
  }{% else
    \setlength{\parindent}{0pt}
    \setlength{\parskip}{6pt plus 2pt minus 1pt}}
}{% if KOMA class
  \KOMAoptions{parskip=half}}
\makeatother
\usepackage{listings}
\newcommand{\passthrough}[1]{#1}
\lstset{defaultdialect=[5.3]Lua}
\lstset{defaultdialect=[x86masm]Assembler}
\usepackage{graphicx}
\makeatletter
\newsavebox\pandoc@box
\newcommand*\pandocbounded[1]{% scales image to fit in text height/width
  \sbox\pandoc@box{#1}%
  \Gscale@div\@tempa{\textheight}{\dimexpr\ht\pandoc@box+\dp\pandoc@box\relax}%
  \Gscale@div\@tempb{\linewidth}{\wd\pandoc@box}%
  \ifdim\@tempb\p@<\@tempa\p@\let\@tempa\@tempb\fi% select the smaller of both
  \ifdim\@tempa\p@<\p@\scalebox{\@tempa}{\usebox\pandoc@box}%
  \else\usebox{\pandoc@box}%
  \fi%
}
% Set default figure placement to htbp
% Make use of float-package and set default placement for figures to H.
% The option H means 'PUT IT HERE' (as  opposed to the standard h option which means 'You may put it here if you like').
\usepackage{float}
\floatplacement{figure}{H}
\makeatother
\setlength{\emergencystretch}{3em} % prevent overfull lines
\providecommand{\tightlist}{%
  \setlength{\itemsep}{0pt}\setlength{\parskip}{0pt}}
\usepackage{bookmark}
\IfFileExists{xurl.sty}{\usepackage{xurl}}{} % add URL line breaks if available
\urlstyle{same}
\definecolor{default-linkcolor}{HTML}{A50000}
\definecolor{default-filecolor}{HTML}{A50000}
\definecolor{default-citecolor}{HTML}{4077C0}
\definecolor{default-urlcolor}{HTML}{4077C0}

\hypersetup{
  hidelinks,
  breaklinks=true,
  pdfcreator={LaTeX via pandoc with the Eisvogel template}}
\author{}
\date{}


%
% for the background color of the title page
%

%
% break urls
%
\PassOptionsToPackage{hyphens}{url}

%
% When using babel or polyglossia with biblatex, loading csquotes is recommended
% to ensure that quoted texts are typeset according to the rules of your main language.
%
\usepackage{csquotes}

%
% captions
%
\definecolor{caption-color}{HTML}{777777}
\usepackage[font={stretch=1.2}, textfont={color=caption-color}, position=top, skip=4mm, labelfont=bf, singlelinecheck=false, justification=raggedright]{caption}
\setcapindent{0em}

%
% blockquote
%
\definecolor{blockquote-border}{RGB}{221,221,221}
\definecolor{blockquote-text}{RGB}{119,119,119}
\usepackage{mdframed}
\newmdenv[rightline=false,bottomline=false,topline=false,linewidth=3pt,linecolor=blockquote-border,skipabove=\parskip]{customblockquote}
\renewenvironment{quote}{\begin{customblockquote}\list{}{\rightmargin=0em\leftmargin=0em}%
\item\relax\color{blockquote-text}\ignorespaces}{\unskip\unskip\endlist\end{customblockquote}}

%
% Source Sans Pro as the default font family
% Source Code Pro for monospace text
%
% 'default' option sets the default
% font family to Source Sans Pro, not \sfdefault.
%
% Note that the font has been officially renamed to `Source Sans 3`, and
% the version provided by the `sourcesanspro` package is slightly outdated.
% You can install the newer version locally and use it, for example, with
% `mainfont: "Source Sans 3"` in the YAML metadata (requires XeTeX or LuaTeX).
%
\ifnum 0\ifxetex 1\fi\ifluatex 1\fi=0 % if pdftex
    \usepackage[default]{sourcesanspro}
  \usepackage{sourcecodepro}
  \else % if not pdftex
    \usepackage[default]{sourcesanspro}
  \usepackage{sourcecodepro}

  % XeLaTeX specific adjustments for straight quotes: https://tex.stackexchange.com/a/354887
  % This issue is already fixed (see https://github.com/silkeh/latex-sourcecodepro/pull/5) but the
  % fix is still unreleased.
  % TODO: Remove this workaround when the new version of sourcecodepro is released on CTAN.
  \ifxetex
    \makeatletter
    \defaultfontfeatures[\ttfamily]
      { Numbers   = \sourcecodepro@figurestyle,
        Scale     = \SourceCodePro@scale,
        Extension = .otf }
    \setmonofont
      [ UprightFont    = *-\sourcecodepro@regstyle,
        ItalicFont     = *-\sourcecodepro@regstyle It,
        BoldFont       = *-\sourcecodepro@boldstyle,
        BoldItalicFont = *-\sourcecodepro@boldstyle It ]
      {SourceCodePro}
    \makeatother
  \fi
  \fi

%
% heading color
%
\definecolor{heading-color}{RGB}{40,40,40}
% By default, the KOMA-Script classes will typeset sectioning headings in
% sans-serif. Use the normal body font for headings.
\addtokomafont{disposition}{\normalfont\color{heading-color}\bfseries}

%
% variables for title, author and date
%
\usepackage{titling}
\title{}
\author{}
\date{}

%
% tables
%

%
% remove paragraph indentation
%
\setlength{\parindent}{0pt}
\setlength{\parskip}{6pt plus 2pt minus 1pt}
\setlength{\emergencystretch}{3em}  % prevent overfull lines

%
%
% Listings
%
%


%
% general listing colors
%
\definecolor{listing-background}{HTML}{F7F7F7}
\definecolor{listing-rule}{HTML}{B3B2B3}
\definecolor{listing-numbers}{HTML}{B3B2B3}
\definecolor{listing-text-color}{HTML}{000000}
\definecolor{listing-keyword}{HTML}{435489}
\definecolor{listing-keyword-2}{HTML}{1284CA} % additional keywords
\definecolor{listing-keyword-3}{HTML}{9137CB} % additional keywords
\definecolor{listing-identifier}{HTML}{435489}
\definecolor{listing-string}{HTML}{00999A}
\definecolor{listing-comment}{HTML}{8E8E8E}

\lstdefinestyle{eisvogel_listing_style}{
  language         = java,
  numbers          = left,
  xleftmargin      = 2.7em,
  framexleftmargin = 2.5em,
  backgroundcolor  = \color{listing-background},
  basicstyle       = \color{listing-text-color}\linespread{1.0}%
                      \lst@ifdisplaystyle%
                      \small%
                      \fi\ttfamily{},
  breaklines       = true,
  frame            = single,
  framesep         = 0.19em,
  rulecolor        = \color{listing-rule},
  frameround       = ffff,
  tabsize          = 4,
  numberstyle      = \color{listing-numbers},
  aboveskip        = 1.0em,
  belowskip        = 0.1em,
  abovecaptionskip = 0em,
  belowcaptionskip = 1.0em,
  keywordstyle     = {\color{listing-keyword}\bfseries},
  keywordstyle     = {[2]\color{listing-keyword-2}\bfseries},
  keywordstyle     = {[3]\color{listing-keyword-3}\bfseries\itshape},
  sensitive        = true,
  identifierstyle  = \color{listing-identifier},
  commentstyle     = \color{listing-comment},
  stringstyle      = \color{listing-string},
  showstringspaces = false,
  escapeinside     = {/*@}{@*/}, % Allow LaTeX inside these special comments
  literate         =
  {á}{{\'a}}1 {é}{{\'e}}1 {í}{{\'i}}1 {ó}{{\'o}}1 {ú}{{\'u}}1
  {Á}{{\'A}}1 {É}{{\'E}}1 {Í}{{\'I}}1 {Ó}{{\'O}}1 {Ú}{{\'U}}1
  {à}{{\`a}}1 {è}{{\`e}}1 {ì}{{\`i}}1 {ò}{{\`o}}1 {ù}{{\`u}}1
  {À}{{\`A}}1 {È}{{\`E}}1 {Ì}{{\`I}}1 {Ò}{{\`O}}1 {Ù}{{\`U}}1
  {ä}{{\"a}}1 {ë}{{\"e}}1 {ï}{{\"i}}1 {ö}{{\"o}}1 {ü}{{\"u}}1
  {Ä}{{\"A}}1 {Ë}{{\"E}}1 {Ï}{{\"I}}1 {Ö}{{\"O}}1 {Ü}{{\"U}}1
  {â}{{\^a}}1 {ê}{{\^e}}1 {î}{{\^i}}1 {ô}{{\^o}}1 {û}{{\^u}}1
  {Â}{{\^A}}1 {Ê}{{\^E}}1 {Î}{{\^I}}1 {Ô}{{\^O}}1 {Û}{{\^U}}1
  {œ}{{\oe}}1 {Œ}{{\OE}}1 {æ}{{\ae}}1 {Æ}{{\AE}}1 {ß}{{\ss}}1
  {ç}{{\c c}}1 {Ç}{{\c C}}1 {ø}{{\o}}1 {å}{{\r a}}1 {Å}{{\r A}}1
  {€}{{\EUR}}1 {£}{{\pounds}}1 {«}{{\guillemotleft}}1
  {»}{{\guillemotright}}1 {ñ}{{\~n}}1 {Ñ}{{\~N}}1 {¿}{{?`}}1
  {…}{{\ldots}}1 {≥}{{>=}}1 {≤}{{<=}}1 {„}{{\glqq}}1 {“}{{\grqq}}1
  {”}{{''}}1
}
\lstset{style=eisvogel_listing_style}

%
% Java (Java SE 12, 2019-06-22)
%
\lstdefinelanguage{Java}{
  morekeywords={
    % normal keywords (without data types)
    abstract,assert,break,case,catch,class,continue,default,
    do,else,enum,exports,extends,final,finally,for,if,implements,
    import,instanceof,interface,module,native,new,package,private,
    protected,public,requires,return,static,strictfp,super,switch,
    synchronized,this,throw,throws,transient,try,volatile,while,
    % var is an identifier
    var
  },
  morekeywords={[2] % data types
    % primitive data types
    boolean,byte,char,double,float,int,long,short,
    % String
    String,
    % primitive wrapper types
    Boolean,Byte,Character,Double,Float,Integer,Long,Short
    % number types
    Number,AtomicInteger,AtomicLong,BigDecimal,BigInteger,DoubleAccumulator,DoubleAdder,LongAccumulator,LongAdder,Short,
    % other
    Object,Void,void
  },
  morekeywords={[3] % literals
    % reserved words for literal values
    null,true,false,
  },
  sensitive,
  morecomment  = [l]//,
  morecomment  = [s]{/*}{*/},
  morecomment  = [s]{/**}{*/},
  morestring   = [b]",
  morestring   = [b]',
}

\lstdefinelanguage{XML}{
  morestring      = [b]",
  moredelim       = [s][\bfseries\color{listing-keyword}]{<}{\ },
  moredelim       = [s][\bfseries\color{listing-keyword}]{</}{>},
  moredelim       = [l][\bfseries\color{listing-keyword}]{/>},
  moredelim       = [l][\bfseries\color{listing-keyword}]{>},
  morecomment     = [s]{<?}{?>},
  morecomment     = [s]{<!--}{-->},
  commentstyle    = \color{listing-comment},
  stringstyle     = \color{listing-string},
  identifierstyle = \color{listing-identifier}
}

%
% header and footer
%
\usepackage[headsepline,footsepline]{scrlayer-scrpage}

\newpairofpagestyles{eisvogel-header-footer}{
  \clearpairofpagestyles
  \ihead*{}
  \chead*{}
  \ohead*{}
  \ifoot*{}
  \cfoot*{}
  \ofoot*{\thepage}
  \addtokomafont{pageheadfoot}{\upshape}
}
\pagestyle{eisvogel-header-footer}



%
% Define watermark
%

\begin{document}

\title{DeFi's First Venture Fund Token - \$FVC}
\author{AKH DIGITAL Team\thanks{info@akh-digital.com}}
\maketitle

\begin{abstract}
\section*{Executive Summary}
The First Venture Capital (\$FVC) token represents a paradigm shift in
decentralised finance: the world's first interest-free,
community-governed token for venture capital and business grants. We're
providing the solution small businesses have been waiting for:
access to capital. Instead of giving away 30\% equity for \$200k,
receive funding through community governance. The community that funds
you becomes your biggest supporters. With a focus on sensible funding
and maintaining complete decentralisation, we aim to provide liquidity
to those that need it, and perpetual profitability.
\end{abstract}

\section{1. Market Problem}\label{market-problem}

Small and medium-sized enterprises (SMEs) face a critical financing gap
in both Web3 and traditional industries. UK early-stage startup funding
has declined to a six-year low, with fewer fundraising rounds despite
rising average deal sizes. Traditional venture capital funding cycles
typically exceed six months, delaying capital deployment and restricting
access to timely capital. This challenge is particularly acute for
FCA-authorised and compliant startups navigating regulatory constraints
whilst seeking growth capital.

The UK government has recognised this financing gap through significant
policy initiatives, including ministerial reports on SME access to
finance and substantial funding packages aimed at supporting small
business growth. Despite these efforts, structural barriers persist:
retail investors remain largely excluded from early-stage opportunities,
DeFi tokens lack clear legal certainty within the UK regulatory
framework, and traditional financing models impose lengthy approval
processes that hinder innovation.

This capital bottleneck restricts innovation, diminishes the
competitiveness of the UK technology sector, and leaves promising
ventures without access to the resources they need to scale. The market
demands a new approach---one that leverages blockchain technology for
transparency and automation, employs tokenisation to make capital
raising more efficient and accessible, and provides decentralised
governance whilst maintaining regulatory compliance. Those who fail to
adapt to this fundamental shift in finance risk being left behind as
tokenised financial infrastructure transforms global markets.

\section{2. Token Solution}\label{token-solution}

\subsection{2.1 Two Simple Functions, Infinite
Possibilities}\label{two-simple-functions-infinite-possibilities}

Web3 companies deposit collateral, get immediate liquidity. No interest,
no banks, no waiting. The token earns fees on every grant.

The community votes to fund promising SMEs
and startups. In return, smart contracts automatically capture a
percentage of their revenues. One successful investment can generate
millions in perpetual income.

\subsection{2.2 The Magic:
Decentralisation}\label{the-magic-decentralisation}

No fund managers. No monopoly on authority. Your wallet. Your vote. Your venture fund.

Token holders ARE the venture capitalists:

\begin{itemize}
\tightlist
\item
  \textbf{Vote} on which SMEs and startups
  to fund
\item
  \textbf{Earn} from every investment
\item
  \textbf{Govern} all token parameters
\item
  \textbf{Own} the entire ecosystem
\end{itemize}

\subsection{2.3 Single-Token System}\label{single-token-system}

AKH DIGITAL introduces a single-token system comprising:

\begin{itemize}
\tightlist
\item{\textbf{\$FVC (Governance Token):}} Tradable, fixed-supply ERC-20 token used for bonding, price discovery, governance rights and revenue distribution.
\item
  \textbf{Bonding Contracts:} Allocate \$FVC at a discount in exchange
  for stable reserves; funds directed to the treasury.
\item
  \textbf{Staking Mechanism:} \$FVC tokens locked for 3--12 months,
  yielding percentage of ecosystem revenue funded from audited SMEs and startups profit repayments.
\item
  \textbf{Quadratic Governance:} Vote power equals square root of staked
  \$FVC to limit concentration.
\item
  \textbf{Compliance Layer:} On-chain KYC via Sumsub verified process;
  whitelist registry enforcing AML and sector eligibility.
\end{itemize}

\section{3. Staking Framework}\label{staking-framework}

The FVC ecosystem implements the Synthetix proportional staking pattern, an audited mechanism deployed across billions in Total Value Locked. Users stake FVC tokens into \texttt{StakingRewards.sol} to earn USDC rewards distributed proportionally based on stake weight and duration.

Crucially, FVC utilises a \textbf{Non-Dilutive Settlement} model: we pay rewards in \textbf{USDC} (revenue) rather than printing new FVC tokens (inflation). This preserves token value by ensuring yield is derived from actual economic output generated by the \textbf{Active Treasury Engine} rather than supply dilution.

The Sustainability Ratio ($S_r$) of the ecosystem is expressed as:

\[
S_r = \frac{\text{Treasury Revenue (USDC)}}{\text{Token Inflation (FVC)}} > 1
\]

In traditional DeFi, $S_r < 1$ (Ponzi dynamics). In FVC, because inflation is minimised and rewards are external (USDC), $S_r$ is designed to trend positive.

The reward distribution follows:

\[
r_i = \frac{s_i}{S} \cdot R \cdot t
\]

$r_i$ represents individual rewards, $s_i$ the user's staked balance, $S$ total staked supply, $R$ the per-second reward rate, and $t$ elapsed time. This system eschews tiered lockups in favour of continuous accrual, maintaining capital efficiency whilst eliminating lock-in friction for participants. Rewards derive from SME revenue-sharing agreements and Real World Asset (RWA) returns, creating sustainable yield backed by productive economic activity.

\begin{enumerate}
    \item{\textbf{Staking Rewards}} - A competitive APY from SME revenue repayments
    \item{\textbf{Priority Access}} - First rights to investment opportunities
    \item{\textbf{Governance Rights}} - Vote on protocol decisions with quadratic weighting
\end{enumerate}
\newpage
\section{4. Tokenomics}\label{token-distribution-mechanics}

\subsection{4.1 Overview}\label{overview-1}

The FVC Ecosystem Token Distribution Mechanics ensure fair, transparent,
and controlled distribution of the \$FVC token supply while maintaining
price stability and preventing market manipulation.

\textbf{Presale (23.0\%)}

\begin{itemize}
\tightlist
\item
  \textbf{Allocation:} 230,000,000 FVC tokens
  (includes seed round, private presale, and all sales before public launch)
\item
  \textbf{Schedule:} Vesting terms subject to negotiation based on
  market conditions
\item
  \textbf{Target:} Raise amount to be determined based on market demand
\item
  \textbf{Note:} Tokens are minted on-demand through bonding contracts.
  Unallocated tokens remain unminted, with the allocation adjusting
  based on actual market demand rather than predetermined distribution.
\end{itemize}

\textbf{Founders, Team \& Partners (25.0\%)}

\begin{itemize}
\tightlist
\item
  \textbf{Allocation:} 250,000,000 FVC tokens
\item
  \textbf{Schedule:} Vesting terms subject to negotiation with institutional investors
\item
  \textbf{Performance Ties:} Milestone-based unlocks with clawbacks
\end{itemize}

\textbf{Marketing \& Community (10.0\%)}

\begin{itemize}
\tightlist
\item
  \textbf{Allocation:} 100,000,000 FVC tokens (provisional)
\item
  \textbf{Governance:} Professional delegates, reviewers, tooling
\item
  \textbf{Airdrops:} Sybil-resistant programmes with retention metrics
\item
  \textbf{Staged Releases:} Milestone-based unlocks with KPIs
\end{itemize}

\textbf{Treasury (35.0\%)}

\begin{itemize}
\tightlist
\item
  \textbf{Allocation:} 350,000,000 FVC tokens
\item
  \textbf{Purpose:} Flexible allocation for ecosystem needs, grants, partnerships, and strategic initiatives
\item
  \textbf{Control:} Multi-signature safe wallet with governance oversight
\item
  \textbf{Note:} Used for operational flexibility and long-term ecosystem development
\end{itemize}

\textbf{Security Reserve (2.0\%)}

\begin{itemize}
\tightlist
\item
  \textbf{Allocation:} 20,000,000 FVC tokens
\item
  \textbf{Purpose:} Emergency security fund for protocol protection and incident response
\item
  \textbf{Governance:} Subject to community governance approval for deployment
\item
  \textbf{Note:} Reserved for security incidents, bug bounties, and protocol defence
\end{itemize}

\textbf{Liquidity Provision (5.0\%)}

\begin{itemize}
\tightlist
\item
  \textbf{Allocation:} 50,000,000 FVC tokens
\item
  \textbf{Initial Pools:} Uniswap V3 (FVC/USDC, FVC/WETH)
\item
  \textbf{Range Management:} Concentrated liquidity for efficiency
\item
  \textbf{Market Quality:} Reduced slippage and better price discovery
\end{itemize}

\textbf{Treasury Operations (USDC)}

\begin{itemize}
\tightlist
\item
  \textbf{Capital Source:} USDC from multi-signature safe wallet for operational expenses
\item
  \textbf{Grant Deployment:} Primary capital pool for SME funding operations
\item
  \textbf{Staking Funding:} APY from grant repayment revenues and liquidation fees
\item
  \textbf{Revenue Allocation:} Revenues allocated to staking yields, reserves, and buy-and-burn operations
\item
  \textbf{Note:} Buy-and-burn mechanism purchases FVC from the market and burns tokens, providing price support without minting new supply
\end{itemize}


\section{5. Technical Architecture}\label{technical-architecture}

FVC is a \textbf{Capital Deployment Operating System}. It is a modular, upgradeable Web3 infrastructure stack focused
on compliant venture funding. It operates on Ethereum-compatible
networks and integrates modern tooling for secure, maintainable
development across contracts, frontend, and infrastructure layers.

\subsection{5.1 Deployed Smart Contracts (MVP)}\label{deployed-smart-contracts}

\textbf{Core Contracts}

\begin{itemize}
\tightlist
\item
  \passthrough{\lstinline!FVC.sol!} --- ERC20 token with 1 billion token cap, role-based minting and burning
  \begin{itemize}
  \tightlist
  \item
    MINTER\_ROLE for Sale contract to mint on-demand
  \item
    BURNER\_ROLE for buy-and-burn operations
  \item
    Standard ERC20Capped implementation with AccessControl
  \end{itemize}
\end{itemize}

\textbf{Token Sale}

\begin{itemize}
\tightlist
\item
  \passthrough{\lstinline!Sale.sol!} --- Fixed-price sale accepting USDC/USDT
  \begin{itemize}
  \tightlist
  \item
    Mints FVC on-demand when purchased
  \item
    Owner-controlled rate, cap, and active status
  \item
    \textbf{Automated Vesting:} Large purchases trigger automatic vesting schedule creation
  \item
    Sends stablecoins to treasury Safe
  \end{itemize}
\end{itemize}

\textbf{Vesting}

\begin{itemize}
\tightlist
\item
  \passthrough{\lstinline!Vesting.sol!} --- Time-based vesting logic
  \begin{itemize}
  \tightlist
  \item
    Enforces linear vesting schedules with optional cliff periods
  \item
    Managed by Gnosis Safe governance or automated via Sale contract
  \item
    Supports revocable schedules for team performance accountability
  \item
    \textbf{Vested Amount Calculation:}
    \[
    V = \begin{cases} 
    0 & \text{if } t < \text{cliff} \\
    \text{total} \times \frac{t - \text{start}}{\text{duration}} & \text{if } t \geq \text{cliff}
    \end{cases}
    \]
  \end{itemize}
\end{itemize}

\textbf{Staking}

\begin{itemize}
\tightlist
\item
  \passthrough{\lstinline!Staking.sol!} --- Synthetix proportional staking pattern
  \begin{itemize}
  \tightlist
  \item
    Stake FVC, earn USDC rewards
  \item
    No lock period, withdraw anytime
  \item
    Owner adds rewards via \passthrough{\lstinline!notifyRewardAmount()!}
  \item
    Battle-tested pattern used by Curve, Uniswap, and SushiSwap
  \end{itemize}
\end{itemize}

\textbf{Testnet Utilities}

\begin{itemize}
\tightlist
\item
  \passthrough{\lstinline!FVCFaucet.sol!} --- Testnet-only token distribution
\item
  \passthrough{\lstinline!MockStable.sol!} --- Test USDC for development
\end{itemize}

\subsection{5.2 Treasury Management (Active Treasury Engine)}\label{treasury-management}

FVC operates an \textbf{Active Treasury Engine} rather than a static vault. While currently secured by a Gnosis Safe Governance Module for the Genesis Phase, the architecture is designed to use \textbf{Yield Adapters} to programmatically route idle capital into low-risk, liquidity-grade strategies (such as Aave or RWA vaults) to generate continuous revenue.

The efficiency of the Active Treasury ($E_t$) is defined as the summation of yield generated across $n$ deployed adapters ($A_i$) minus operational latency ($L_o$):

\[
E_t = \sum_{i=1}^{n} (A_i \cdot r_i) - L_o
\]

Where $r_i$ is the risk-adjusted return rate of adapter $i$. Unlike static treasuries where $E_t \approx 0$, FVC maximises $E_t$ through algorithmic rebalancing.

\textbf{Multi-Signature Safe Wallet (3-of-5)}

\begin{itemize}
\tightlist
\item
  Holds USDC reserves and manages \textbf{Algorithmic Risk Management} parameters
\item
  Manages FVC token minting permissions
\item
  Controls TokenSale and Staking contracts
\end{itemize}

\textbf{Safe Transaction Flow}

\begin{enumerate}
\def\labelenumi{\arabic{enumi}.}
\tightlist
\item
  Proposal created by signer
\item
  Requires 3 of 5 signatures to execute
\item
  Executes on-chain via Safe proxy contract
\end{enumerate}

\textbf{Common Operations}

\begin{itemize}
\tightlist
\item
  Grant TokenSale MINTER\_ROLE: \passthrough{\lstinline!fvc.grantRole(MINTER\_ROLE, saleAddress)!}
\item
  Add staking rewards: \passthrough{\lstinline!staking.notifyRewardAmount(usdcAmount)!}
\item
  Mint team tokens: \passthrough{\lstinline!fvc.mint(teamWallet, amount)!}
\item
  Buy-and-burn: Swap USDC for FVC on DEX, then \passthrough{\lstinline!fvc.burn(amount)!}
\item
  Update sale parameters: \passthrough{\lstinline!sale.setRate()!}, \passthrough{\lstinline!sale.setCap()!}, \passthrough{\lstinline!sale.setActive()!}
\end{itemize}

\textbf{Safe Integration}

\begin{itemize}
\tightlist
\item
  ABI interfaces allow Safe to call any contract function
\item
  No custom integration contracts needed
\item
  All operations transparent and auditable on-chain
\end{itemize}

\subsection{5.3 Contract ABIs \& External Integration}\label{contract-abis-external-integration}

\textbf{Application Binary Interface (ABI)}

\begin{itemize}
\tightlist
\item
  JSON specification of contract functions and events
\item
  Enables frontend, Safe wallet, and external services to interact with contracts
\item
  Generated automatically during compilation
\end{itemize}

\textbf{Key Integrations}

\begin{itemize}
\tightlist
\item
  \textbf{Frontend (dApp) → Contracts:} Uses ethers.js with TypeChain for type-safe calls
\item
  \textbf{Gnosis Safe → Contracts:} Safe UI reads contract ABIs, displays function parameters for signers, constructs transaction calldata, executes after threshold signatures collected
\item
  \textbf{Backend Services → Contracts:} Event monitoring for TokenSale purchases, staking reward distribution automation, revenue tracking for buy-and-burn
\end{itemize}

\textbf{Event Emissions}

\begin{itemize}
\tightlist
\item
  All state changes emit events for off-chain tracking
\item
  Frontend and backend listen to events for real-time updates
\item
  Examples: TokensPurchased, Staked, RewardPaid
\end{itemize}

\subsection{5.4 Planned Contracts (Phase 2+)}\label{planned-contracts}

\textbf{Governance Infrastructure}

\begin{itemize}
\tightlist
\item
  On-chain voting for protocol parameters
\item
  Proposal creation and execution
\item
  Timelock for security delays
\item
  Emergency pause mechanisms
\end{itemize}

\textbf{Advanced Compliance}

\begin{itemize}
\tightlist
\item
  KYC/AML verification integration
\item
  On-chain whitelist management
\item
  Regulatory reporting tools
\end{itemize}

\textbf{Treasury Automation}

\begin{itemize}
\tightlist
\item
  Automated revenue distribution
\item
  Grant milestone releases
\item
  Vesting schedules for team and investors
\end{itemize}

\textbf{Additional Staking}

\begin{itemize}
\tightlist
\item
  Tiered lock-up periods with APY multipliers
\item
  LP token staking
\item
  NFT staking for governance weight
\end{itemize}

\textbf{Note:} Phase 2 contracts will be developed, audited, and deployed based on ecosystem needs and governance decisions

%%%
%\begin{figure}
%\centering
%\pandocbounded{\includegraphics[keepaspectratio]{images/ecosystem.png}}
%\caption{FVC Ecosystem Architecture}
%\end{figure}
%%%
\subsection{5.5 Security Measures}\label{security-measures}

\begin{itemize}
\tightlist
\item
  \textbf{Multi-Signature Security}

  \begin{itemize}
  \tightlist
  \item
    \textbf{Multi-signature wallets} requiring 3 out of 5 authorised
    signers (founders) to approve any transaction
  \item
    \textbf{Treasury operations} protected by Gnosis Safe implementation
    with role-based access control
  \item
    \textbf{Emergency functions} require consensus among guardian
    signers
  \end{itemize}
\item
  \textbf{Timelock Implementation}

  \begin{itemize}
  \tightlist
  \item
    \textbf{Timelock delay} of 48 hours between governance proposal
    approval and execution
  \item
    \textbf{Gives token holders time to react} if they disagree with
    decisions or identify security concerns
  \item
    \textbf{Variable delays} based on operation criticality (7 days for
    critical operations)
  \item
    \textbf{Maximum delay cap} of 30 days for complex multi-step
    operations
  \end{itemize}
\item
  \textbf{Guardian Powers}

  \begin{itemize}
  \tightlist
  \item
    \textbf{Special emergency functions} that can immediately stop
    system functions during crises
  \item
    \textbf{Can pause/unpause smart contracts} to prevent exploits
  \item
    \textbf{Funds cannot be moved} during emergency pauses
  \item
    \textbf{Governance can be bypassed} during emergency situations
  \item
    \textbf{3-of-5 multisig controls this mechanism}
  \end{itemize}
\end{itemize}

\subsection{5.3 Upgrade Mechanisms}\label{upgrade-mechanisms}

\begin{itemize}
\tightlist
\item
  \textbf{UUPS Proxy Implementation}
\end{itemize}

The ecosystem implements OpenZeppelin's UUPS (Universal Upgradeable Proxy
Standard), a smart contract design pattern that allows developers to
upgrade contract logic while preserving the contract's address, state
and user balances. It works by delegating function calls to a nested
module within a contract, while FVC.sol contains a permanent address
users interact with, the actual logic is delegated to another
implementation which is upgradeable, changing nothing from the user's
perspective.

\subsection{5.4 Integration Architecture: How You Connect with Borrower
Systems}\label{integration-architecture-how-you-connect-with-borrower-systems}

\begin{itemize}
\tightlist
\item
  \textbf{External API Integration}
\end{itemize}

The system connects with established financial and legal infrastructure
through secure API integrations. Sumsub provides KYC/AML verification
services with real-time status updates. DocuSign enables automated legal
document management and digital signature collection. Banking APIs
facilitate traditional payment processing and account verification for
SME borrowers.

All external integrations implement proper authentication, rate
limiting, and error handling to ensure reliable operation. Fallback
mechanisms and manual processing capabilities provide continuity during
external service disruptions.

\begin{itemize}
\tightlist
\item
  \textbf{Oracle Integration Framework}
\end{itemize}

The ecosystem integrates with multiple oracle systems for critical data
feeds. Chainlink price oracles provide real-time FVC/USD pricing for
bonding calculations and treasury management. Custom KYC compliance
oracles connect with external identity verification providers, enabling
automated eligibility verification while maintaining user privacy.

Oracle integration includes redundancy and validation mechanisms to
prevent manipulation or single points of failure. Price feeds implement
staleness checks and circuit breakers to halt operations during oracle
failures. KYC oracles utilise cryptographic attestations and
multi-provider verification for enhanced security.

\begin{itemize}
\tightlist
\item
  \textbf{Borrower System Integration}
\end{itemize}

The ecosystem connects with SME borrowers through standardised
revenue-sharing contract interfaces, enabling automated payment
processing and transparent revenue distribution. Borrowers deploy
ecosystem-compliant contracts that handle payment schedules, late fee
calculations, and automatic revenue forwarding to the ecosystem treasury.

The integration system supports multiple payment methods including
direct cryptocurrency transfers, traditional bank integration through
Plaid/Yodlee APIs, and manual payment processing for complex
arrangements. All payments are tracked on-chain with immutable records
for audit and compliance purposes.

\subsection{Data Architecture: On-Chain vs Off-Chain Data
Strategy}\label{data-architecture-on-chain-vs-off-chain-data-strategy}

\begin{itemize}
\tightlist
\item
  \textbf{On-Chain Data Strategy}
\end{itemize}

Critical ecosystem data including token balances, governance votes,
vesting schedules, and parameter configurations are stored directly on
the blockchain, ensuring immutability, transparency, and decentralised
access. Smart contracts emit comprehensive events for all state changes,
enabling efficient indexing and real-time monitoring.

On-chain storage optimisation reduces gas costs through packed data
structures, efficient storage layouts, and event-based rather than
storage-based logging where appropriate. The system prioritises data
integrity and accessibility while managing storage costs through careful
architectural decisions.

\begin{itemize}
\tightlist
\item
  \textbf{Off-Chain Data Management}
\end{itemize}

Privacy-sensitive information including KYC documents, legal agreements,
and business plans are stored off-chain using the InterPlanetary File
System (IPFS) for decentralised document storage. Documents are
encrypted where necessary and referenced on-chain through immutable
content hashes, ensuring integrity while maintaining privacy compliance.

The off-chain storage strategy includes redundancy through multiple IPFS
pinning services, automated backup procedures, and document lifecycle
management. Access controls ensure that sensitive documents are only
available to authorised parties while maintaining transparency for
public information.

\subsection{5.6 Tooling Stack}\label{tooling-stack}

\begin{itemize}
\tightlist
\item
  \textbf{Hardhat:} Full-stack Ethereum development environment for
  compiling, testing, deploying, and upgrading Solidity contracts.
\item
  \textbf{OpenZeppelin:} Secure libraries used for implementing ERC20,
  access control, and upgradeability via proxy patterns.
\item
  \textbf{Node.js \& npm:} JavaScript runtime (v18+ LTS recommended)
  with npm for dependency management and script execution.
\item
  \textbf{TypeChain:} Generates type-safe TypeScript bindings from
  contract ABIs for use in both backend and frontend layers.
\item
  \textbf{wagmi \& viem:} TypeScript-native libraries for wallet
  connections and contract interactions in React-based dApps.
\end{itemize}

\section{6. Dual-Channel Presale Architecture}\label{dual-channel-presale-architecture}

\subsection{6.1 Presale Overview}\label{presale-overview}

\textbf{Total Allocation:} 230,000,000 tokens (23.0\% of total supply) distributed through seed round, private presale, and all sales before public launch.

\textbf{Vesting Terms:} Subject to negotiation with institutional investors based on market conditions and investor requirements.

\subsection{Community Channel (Bonding Mechanism)}\label{community-channel-bonding-mechanism}

The pre-sale phase incorporates a bonding mechanism with milestone-based pricing and controlled token distribution. This approach ensures controlled access whilst maintaining transparent price discovery.

The pricing is programmed as a reverse Dutch auction where the price starts low and gets higher to reward early clients. The only caveat is the prices have milestones rather than a gradual increase.

\subsection{6.3 Private Channel (OTC Swap)}\label{private-channel-otc-swap}

Parallel private channel for accredited investors featuring fixed pricing, allowlist verification, and custom allocation limits. Provides negotiated terms for institutional participants.

\section{7. Grants System}\label{grants-system}

\subsection{7.1 Dual Grant System}\label{dual-grant-system}

FVC offers two grant categories designed to serve different
SME and startup stages and capital needs.
All grants are interest-free and structured to align incentives between
the ecosystem and funded projects.

\subsubsection{Small Grants}\label{small-grants}

Small grants provide fast, accessible funding for early-stage projects
through a founder-friendly structure. Borrowers deposit minimal
collateral in yield-bearing assets. Projects share a small percentage of
revenue over a fixed period---no interest, no predetermined returns,
just aligned incentives. Fast-track vetting ensures quick deployment for
hackathon winners, MVP development, and bridge funding. Collateral only
serves as fraud protection, never liquidated for revenue shortfalls.
This model beats traditional lending terms while generating sustainable
ecosystem returns through actual business performance rather than
interest.

\textbf{Structure:}

\begin{itemize}
\tightlist
\item
  \textbf{Amount:} \$25,000 - \$100,000
\item
  \textbf{Collateral Requirement:} 10-30\% put down in stable assets:
  BTC, ETH, USDC et cetera, helping to protect staked funds.
\item
  \textbf{Revenue Share:} A percentage of gross revenue for 18-24 months
\item
  \textbf{Approval:} Fast-track vetting (72 - 168 hours)
\item
  \textbf{Liquidation:} Only triggered by fraud, misuse of funds, or
  project abandonment
\end{itemize}

\subsubsection{7.1.2 Large Strategic Grants}\label{large-strategic-grants}

Large grants fund established firms through two models: revenue share
agreements that preserve founder equity, or direct token/equity
investments for strategic partnerships. Projects undergo community
vetting where FVC token holders evaluate team, technology, and market
fit. Selected projects access capital plus FVC's ecosystem of advisors
and community supporters. Revenue shares run for a fixed term or
perpetually, generating sustainable yield for stakers without dilution.
Token/equity returns flow to the treasury and stakers constantly and
upon exit. This dual approach supports projects whether they're seeking
non-dilutive capital or established projects offering strategic
allocations.

\textbf{Type A: Revenue Share Agreement}

\textbf{Structure:}

\begin{itemize}
\tightlist
\item
  \textbf{Amount:} \$100,000 - \$500,000+
\item
  \textbf{Duration:} Perpetual or a fixed term
\item
  \textbf{Revenue Share:} A percentage of monthly gross revenue
\item
  \textbf{Approval:} Community vote after comprehensive vetting
\item
  \textbf{Collateral:} None required
\end{itemize}

\textbf{Type B: Equity/Token Investment}

\textbf{Structure:}

\begin{itemize}
\tightlist
\item
  \textbf{Amount:} \$100,000 - \$500,000+
\item
  \textbf{Investment Type:} Direct equity or token allocation
\item
  \textbf{Ownership:} A fixed percentage depending on stage and amount
\item
  \textbf{Approval:} Community vote with enhanced due diligence and the
  same comprehensive vetting
\item
  \textbf{Rights:} Standard investor protections (pro-rata, information
  rights)
\end{itemize}

SMEs and funded firms remit an agreed share of revenues at fixed time periods. Treasury allocates these revenues to: staking yields, FVC reserves and to buy and burn FVC.

\subsubsection{7.1.3 The Vetting Advantage}\label{the-vetting-advantage}

Unlike pure DAOs where anyone can propose anything, FVC employs a
specialised team that filters applications first:

\begin{enumerate}
\def\labelenumi{\arabic{enumi}.}
\tightlist
\item
  SMEs and startups Apply → Professional
  pitch deck, revenue model, team background
\item
  Expert Review → Our team eliminates obvious failures (90\% filtered)
\item
  Due Diligence → Deep dive on the top 10\%
\item
  Community Vote → Only pre-vetted, high-quality projects reach token
  holders
\end{enumerate}

Token holders vote on winners, rather than wading through garbage.
Professional filtering + community wisdom = superior returns.

\subsubsection{How Funding Works}\label{how-funding-works}

Application → Vetting (1 week) → Community Vote (3 days) → Instant
Funding

\textbf{Smart Contract Deployment:}

\begin{itemize}
\tightlist
\item
  \textbf{Approved projects get custom revenue-sharing contracts}
\item
  \textbf{A percentage of all revenues flow back automatically}
\item
  \textbf{No trust required - code enforces the agreement}
\item
  \textbf{Forever returns from one-time investments}
\end{itemize}

\textbf{What Projects Receive:}

\begin{itemize}
\tightlist
\item
  \textbf{Stable assets directly to their wallets (e.g.~ETH, BTC, USDC)}
\item
  \textbf{Marketing support from our community}
\item
  \textbf{Follow-on funding priority}
\end{itemize}

Professional vetting eliminates 90\% of noise. Community governance
selects from curated excellence. Every funded project has passed both
institutional-grade diligence AND community approval. Higher success
rates, better returns, happier token holders.

\subsection{Compliance and
Verification}\label{compliance-and-verification}

FVC implements a comprehensive revenue verification system that
applies uniformly to all grant recipients, ensuring transparency and
accountability across the portfolio. By combining on-chain automation
with off-chain attestation, FVC maintains accurate revenue
tracking while minimising trust assumptions. This hybrid approach
enables verification of both crypto-native and traditional revenue
streams, creating a robust framework for sustainable revenue sharing.

Whitelist mechanism ties wallet address to a zero-knowledge proof
generated via Sumsub KYC process, the approved KYC and identity
attestation provider. Periodic identity verification ensures AML
compliance and sector-specific eligibility. No alternative providers
accepted without governance proposal and approval. Quarterly legal
reviews mandated. Registry excludes gambling, meme coins, adult content,
and unlicensed finance.

FVC's verification team, led by our CFA Accountant specialising
in revenue analysis, actively monitors for revenue manipulation tactics
including restructuring schemes, related-party transactions, and
artificial revenue deferrals. This professional oversight ensures
revenue sharing captures genuine business performance, with any attempts
to game the system triggering immediate investigation and potential
default proceedings.

\subsubsection{7.2.1 Enforcement and
Punishments}\label{enforcement-and-punishments}

FVC enforces revenue reporting through progressive penalties
designed to maintain compliance while allowing for operational
realities. Monthly reports receive a 7-day grace period for late
submission, after which the issue is escalated to the compliance team.
Missing filings result in immediate suspension from ecosystem
benefits including additional funding, community support, and governance
participation. Material misrepresentation or persistent non-compliance
triggers formal default proceedings, with smart contracts automatically
calculating accrued obligations and initiating recovery processes. All
enforcement actions are transparent, with status updates published to
maintain accountability while giving struggling projects opportunity to
remedy before severe consequences.

\subsection{Vetting Process}\label{vetting-process}

Grant applications undergo two-stage evaluation to ensure quality while
preserving decentralised governance. FVC's specialised due diligence
team---comprising chartered accountants, technical analysts, and
industry veterans---assesses each project's financial health, revenue
authenticity, technical feasibility, and market viability. Only projects
meeting strict operational and financial criteria advance to community
voting, accompanied by comprehensive vetting reports. This combination
of professional diligence and community governance ensures informed
capital allocation decisions while significantly reducing default risk
across the portfolio.

\section{8. Governance}\label{governance-section}

\subsection{8.1 Overview}\label{overview}

The FVC ecosystem contains a sophisticated, multi-layered governance
system designed to balance community participation with operational
efficiency. Our governance framework draws inspiration from successful
DeFi protocols while introducing innovative mechanisms specifically
tailored for SME funding operations.

\subsection{8.2 Core Governance
Architecture}\label{core-governance-architecture}

\textbf{Governance Token}: \$FVC serves as the primary governance token
with a fixed supply of 1 billion tokens. Each \$FVC token represents
voting power in ecosystem decisions, with additional mechanisms to
prevent concentration of power.

\textbf{Quadratic Voting Implementation}: FVC employs quadratic
voting to ensure democratic participation while mitigating whale
dominance. Voting power is calculated as the square root of
staked/locked tokens, creating a more equitable distribution of
influence. The Cost of Influence ($C_I$) for a desired Vote Weight ($W_v$) increases exponentially:

\[
C_I = (W_v)^2
\]

Therefore, to double influence ($2W_v$), an entity must stake four times the capital ($4C_I$), making governance capture economically inefficient.

\begin{center}
\text{Voting Power} = \sqrt{\text{Staked Tokens}}
\end{center}

{Voting Mechanisms: Quorum Requirements, Vote
Delegation}\label{voting-mechanisms-quorum-requirements-vote-delegation}

\begin{itemize}
\tightlist
\item
  \textbf{Quorum Requirements:} 10\% of current veFVC supply must
  participate for vote validity. This ensures meaningful community engagement; prevents low participation rates and allows for dynamic adjustment based on ecosystem maturity.

  \textbf{Vote Delegation System:}

  \begin{itemize}
  \tightlist
  \item
    \textbf{Delegate Selection:} Holders can delegate votes to trusted
    community representatives
  \item
    \textbf{Transparency:} Public profiles and voting records for all
    delegates
  \item
    \textbf{Accountability:} Performance tracking and community
    oversight
  \item
    \textbf{Institutional Support:} Professional governance management
    for large holders
  \end{itemize}
\item
  \textbf{Voting Power Distribution:}

  \begin{itemize}
  \tightlist
  \item
    \textbf{Quadratic Voting:} Natural reduction of whale influence
  \item
    \textbf{Staking Requirements:} Voting power tied to staked/locked
    tokens
  \item
    \textbf{Anti-Flash Loan:} 1-day delay prevents borrowed-power
    manipulation
  \end{itemize}
\end{itemize}

\subsection{8.3 Emergency Procedures: Circuit Breakers, Emergency
Pause}\label{emergency-procedures-circuit-breakers-emergency-pause}

\begin{itemize}
\tightlist
\item
  \textbf{Pause Guardian System:}

  \begin{itemize}
  \tightlist
  \item
    \textbf{Multi-signature wallet (3 of 5 signers) with emergency pause
    capabilities}
  \item
    \textbf{Scope:} Can pause bonding, staking, and treasury operations
  \item
    \textbf{Limitations:} Cannot access funds or bypass governance
    mechanisms
  \item
    \textbf{Cooldown:} 24-hour minimum between emergency pauses
  \item
    \textbf{Maximum Duration:} 7-day pause limit with governance review
  \end{itemize}
\item
  \textbf{Emergency Response Protocol:}

  \begin{itemize}
  \tightlist
  \item
    \textbf{Immediate Activation:} Pause capability for security
    incidents
  \item
    \textbf{Community Notification:} Alert within 1 hour of activation
  \item
    \textbf{Governance Review:} Community assessment within 24 hours
  \item
    \textbf{Recovery Process:} Automatic unpause after cooldown unless
    extended
  \end{itemize}
\item
  \textbf{On-Chain Circuit Breaker:}

  \begin{itemize}
  \tightlist
  \item
    An automated \textbf{Algorithmic Risk Management} tool that mechanically suspends protocol operations during extreme market volatility to protect asset backing.
  \item
    \textbf{Treasury Limits:} Daily and single transfer caps prevent
    excessive outflows
  \item
    \textbf{Parameter Bounds:} Hard limits on critical ecosystem
    parameters
  \item
    \textbf{Emergency Fund:} Dedicated resources for crisis response
  \end{itemize}
\end{itemize}

\subsection{8.4 Governance Attack Prevention: Anti-Manipulation
Measures}\label{governance-attack-prevention-anti-manipulation-measures}

\begin{itemize}
\tightlist
\item
  \textbf{Voting Power Protection:}

  \begin{itemize}
  \tightlist
  \item
    \textbf{Snapshot Mechanism:} Voting power measured at fixed block
    height
  \item
    \textbf{Anti-Flash Loan:} 1-day delay prevents borrowed-power
    manipulation
  \item
    \textbf{Quadratic Voting:} Naturally reduces whale influence
  \item
    \textbf{Staking Requirements:} Voting power tied to locked tokens
  \end{itemize}
\item
  \textbf{Identity Verification:}

  \begin{itemize}
  \tightlist
  \item
    \textbf{KYC Integration:} Sumsub implementation for regulatory
    compliance
  \item
    \textbf{Whitelist Management:} Controlled access to bonding and
    governance
  \item
    \textbf{Delegate Disclosures:} Public profiles for all voting
    representatives
  \item
    \textbf{Reputation System:} Track record of successful proposals and
    voting
  \end{itemize}
\item
  \textbf{Technical Safeguards:}

  \begin{itemize}
  \tightlist
  \item
    \textbf{Reentrancy Protection:} Standard OpenZeppelin security
    patterns
  \item
    \textbf{Access Controls:} Role-based permissions for administrative
    functions
  \item
    \textbf{Event Logging:} Complete audit trail for all governance
    actions
  \item
    \textbf{Upgrade Controls:} All contract modifications require
    governance approval
  \end{itemize}
\end{itemize}

\section{9. Distribution
Infrastructure}\label{distribution-infrastructure}

\textbf{Smart Contract System}

\begin{itemize}
\tightlist
\item
  \textbf{VestingVault Contract:} Automated token distribution and
  vesting management
\item
  \textbf{Multi-Signature Control:} Treasury and guardian oversight
\item
  \textbf{Emergency Pause:} Ability to halt distribution in emergencies
\item
  \textbf{Audit Trail:} Complete on-chain record of all distributions
\end{itemize}

\textbf{Multi-Signature Treasury}

\begin{itemize}
\tightlist
\item
  \textbf{Gnosis Safe:} 3-of-5 multi-signature wallet
\item
  \textbf{Founder Signers:} 3 founding team members
\item
  \textbf{Advisor Signers:} 2 strategic advisors
\item
  \textbf{Transaction Requirements:} Minimum 3 signatures
\end{itemize}

\section{10. Market Making Strategy}\label{market-making-strategy}

\textbf{Initial Liquidity Provision}

\begin{itemize}
\tightlist
\item
  \textbf{DEX Pools:} Uniswap V3 (FVC/USDC), Uniswap V3 (FVC/WETH)
\item
  \textbf{Initial Liquidity:} 50,000,000 FVC + equivalent USDC/ETH
\item
  \textbf{Range Management:} Concentrated liquidity for capital
  efficiency
\end{itemize}

\textbf{Liquidity Mining Incentives}

\begin{itemize}
\tightlist
\item
  \textbf{Base Rewards:} Yield for FVC/USDC liquidity providers
\item
  \textbf{Tiered Structure:} Higher rewards for longer-term provision
\item
  \textbf{Performance Metrics:} Volume and retention-based bonuses \#
  14. Risk Management \& Exit Strategy
\end{itemize}

\section{11. Risks and Mitigations}\label{risks-and-mitigations}

\subsection{11.1 Purpose}\label{purpose}

The purpose of this Risk Management Policy is to establish a structured
framework for identifying, assessing, mitigating, and monitoring risks
associated with FVC providing funding to SMEs and startups in the cryptocurrency and blockchain
industry. This policy ensures that AKH Digital Ltd.~operates in line
with industry best practices, regulatory expectations, and
Shariah-compliant financing principles.

\subsection{11.2 Due Diligence}\label{due-diligence}

The FVC Team evaluates grant applicants across multiple risk dimensions
before advancing projects to community vote. Our systematic approach
ensures only financially stable, technically sound, and credibly led
projects receive funding, protecting ecosystem resources while
maintaining high quality standards.

\textbf{Core Requirements:}

\begin{itemize}
\tightlist
\item
  \textbf{Financial Health:} Minimum 6-12 months runway demonstrated
  through bank statements and financial projections
\item
  \textbf{Team Verification:} Comprehensive background checks on all
  founders and directors including KYC, credit history, and reputational
  assessment
\item
  \textbf{Technical Assessment:} Code review, architecture evaluation,
  and security audit status for technical projects
\item
  \textbf{Risk Scoring:} Quantitative matrix evaluating financial
  stability, team credibility, technical soundness, and regulatory
  compliance
\item
  \textbf{Minimum Threshold:} Projects must achieve predetermined risk
  score across all categories to proceed to community voting
\end{itemize}

Projects failing any single criterion are rejected without exception,
ensuring consistent standards across the portfolio.

\subsection{11.3 Collateral and Security}\label{collateral-and-security}

FVC mitigates default risk through a robust collateral framework for
small grants, requiring borrowers to secure funding with digital assets
held in ecosystem-controlled smart contracts. This system provides
capital protection while maintaining transparent, automated enforcement
mechanisms that eliminate subjective decision-making in default
scenarios.

\textbf{Collateral Framework:}

\begin{itemize}
\tightlist
\item
  \textbf{Accepted Assets:} BTC, ETH, and stablecoins (USDC, USDT) held
  in multi-signature escrow contracts
\item
  \textbf{Collateralisation Ratio:} Tiered system depending on grant
  size and risk level. 10-30\% of grant value, with borrowers retaining
  staking yields
\item
  \textbf{Monitoring System:} Real-time mark-to-market valuation with
  automated alerts
\item
  \textbf{Margin Calls:} Triggered when collateral value falls below
  25\%, requiring top-up within 48 hours
\item
  \textbf{Liquidation Protocol:} Automatic liquidation upon three missed
  payments or collateral dropping below threshold
\item
  \textbf{Recovery Process:} Collateral sold to recover grant principal,
  excess returned to borrower
\end{itemize}

\subsection{11.4 Revenue}\label{revenue}

FVC structures repayments through revenue-sharing agreements rather than
interest-bearing loans, aligning ecosystem success with borrower growth
while eliminating fixed payment pressure during low-revenue periods.
Smart contract automation and transparent reporting requirements ensure
consistent revenue capture without relying on borrower goodwill.

\textbf{Revenue Structure:}

\begin{itemize}
\tightlist
\item
  \textbf{Sharing Model:} Fixed percentage of gross revenues rather than
  interest payments
\item
  \textbf{Payment Automation:} Smart contracts automatically route
  revenue shares from on-chain income
\item
  \textbf{Web3 Reporting:} Direct on-chain revenue verification through
  ecosystem integrations
\item
  \textbf{Traditional Business:} Monthly bank statement reconciliation;
  quarterly financial statements verified by chartered accountants as
  well as full annual audits
\item
  \textbf{Enforcement:} Three missed payments trigger default
  proceedings and collateral liquidation where applicable
\end{itemize}

This approach ensures sustainable repayment aligned with actual business
performance rather than arbitrary payment schedules.

\subsection{11.5 Treasury Concentration}\label{treasury-concentration}

FVC maintains strict concentration limits to ensure sustainable revenue
generation and minimise single-point-of-failure risks across the grant
portfolio. The treasury deploys capital across diverse sectors and
geographies, preventing overexposure to any single project, sector, or
regulatory jurisdiction.

\textbf{Concentration Limits:}

\begin{itemize}
\tightlist
\item
  \textbf{Per-Project Cap:} Capped amount of available treasury allocated
  to any single grant recipient. Subject to change
\item
  \textbf{Sector Diversification:} Grants distributed equally across
  varied sectors on-chain and off
\item
\textbf{Geographic Distribution:} Projects spanning multiple
jurisdictions to reduce regulatory concentration
\item
  \textbf{Revenue Source Variety:} Balance between established firms
  and emerging projects spanning Web2, Web3 and real world firms.
\item
  \textbf{Staged Deployment:} Initial grants limited to a smaller
  percentage of treasury, increasing only after successful repayment
  history
\end{itemize}

This disciplined approach to capital deployment ensures the ecosystem
maintains diverse revenue streams while protecting treasury resources
from concentrated default risk.

\subsection{11.6 Default and Insolvency}\label{default-and-insolvency}

When funded projects become insolvent or cease operations, FVC employs
multiple recovery mechanisms to minimise losses and protect treasury
resources. Our layered approach combines immediate collateral recovery
with longer-term asset claims, ensuring maximum value recovery while
maintaining transparent loss accounting for token holders.

\textbf{Recovery Mechanisms:}

\begin{itemize}
\tightlist
\item
  \textbf{Collateral Liquidation:} Immediate conversion of escrowed
  collateral to recover principal
\item
  \textbf{Asset Recovery Rights:} Contractual claims on intellectual
  property, token reserves, or equity stakes in parent entities
\item
  \textbf{Loss Buffer Reserve:} A percentage of treasury maintained in
  stablecoins to absorb expected defaults and aid in mitigating risk
\item
  \textbf{Transparent Write-Offs:} Losses recorded according to
  established accounting standards with public reporting
\end{itemize}

Default proceedings prioritise speed of recovery and capital
preservation, with recovered funds redistributed to active grants rather
than individual compensation.

\subsection{11.7 Counterparty and Compliance
Risk}\label{counterparty-and-compliance-risk}

FVC maintains strict compliance standards through comprehensive
KYC/AML procedures and continuous monitoring of all grant recipients.
This proactive approach prevents regulatory violations, protects the
ecosystem from sanctions exposure, and ensures all funded projects meet
international compliance standards.

\textbf{Compliance Framework:}

\begin{itemize}
\tightlist
\item
  \textbf{Identity Verification:} Mandatory KYC/AML checks including
  Ultimate Beneficial Owner (UBO) identification
\item
  \textbf{Blockchain Analytics:} Continuous monitoring transaction
  screening
\item
  \textbf{Ongoing Surveillance:} Real-time alerts for suspicious
  activity or regulatory changes
\item
  \textbf{Exclusion Criteria:} Automatic rejection of projects with
  sanctions exposure or regulatory violations
\item
  \textbf{Quarterly Reviews:} Regular compliance audits of all active
  grant recipients
\item
  \textbf{Geographic Restrictions:} Prohibited jurisdictions list
  maintained and updated monthly
\end{itemize}

\subsection{11.8 Regulation, Jurisdiction and Shariah
Compliancy}\label{regulation-and-shariah-compliancy}

FVC operates within established regulatory frameworks while
maintaining ethical finance principles through its revenue-sharing
model. The ecosystem's structure deliberately avoids interest-based
returns, aligning with both modern regulatory requirements and Islamic
finance principles through genuine risk-sharing arrangements.

\textbf{Jurisdictional Strategy:}

The project is legally structured in the British Virgin Islands (BVI)
based on specific legal counsel to mitigate Collective Investment
Scheme (CIS) risks within the United Kingdom. This offshore structuring
provides a robust legal foundation for the initial deployment phases.

\textbf{UK Regulatory Roadmap (2028-2029):}

FVC maintains active, provisional engagement with UK government
officials, including the Office of the Business Secretary. The
strategic roadmap targets the establishment of a fully regulated
exchange within the UK jurisdiction by 2028/29, contingent upon
financial stability milestones and the evolving UK crypto-asset
regulatory framework. This long-term vision aligns with growing UK
appetite for crypto-asset innovation and job creation.

\textbf{Compliance Measures:}

\begin{itemize}
\tightlist
\item
  \textbf{Regulatory Alignment:} Strategic BVI incorporation with
  roadmap for UK regulatory integration
\item
  \textbf{Revenue-Sharing Structure:} Interest-free model based on
  actual business performance rather than predetermined returns
\item
  \textbf{Shariah Principles:} Independent Shariah advisory panel to
  consistently ensure regulations are followed
\item
  \textbf{Legal Reviews:} Annual legal assessments to ensure continued
  compliance
\item
  \textbf{Transparent Reporting:} Public compliance updates and
  regulatory correspondence
\end{itemize}

This dual approach ensures broad accessibility while maintaining
regulatory compliance across jurisdictions.

\section{12. Exit \& Recovery Strategy}\label{exit-recovery-strategy}

\begin{itemize}
\tightlist
\item
  \textbf{Revenue-Based Exit:} Companies repay via revenue share until
  fixed term is reached or perpetually
\item
  \textbf{Equity Conversion:} In certain cases, revenue-sharing
  agreements may convert into equity
\item
  \textbf{Collateral Exit:} Default triggers automatic collateral
  liquidation
\item
  \textbf{Token Exit:} Token allocations received as part of agreements
  either liquidated or held as long-term treasury assets
\end{itemize}

\end{document}
